\section*{Overview}

Dinic is the max-flow algorithm I expect to implement in contests unless the graph is tiny or the reduction is very
specialized. It is general, predictable, and usually fast enough for the standard modeling problems built on source,
sink, capacities, and cut arguments.

\section*{When to Use It}

Use Dinic when:

\begin{itemize}
  \item the problem can be modeled as moving as much flow as possible through a capacitated network,
  \item you need edge-disjoint or vertex-disjoint path style reasoning,
  \item the graph is too large for a naive augmenting-path implementation.
\end{itemize}

Typical reductions include bipartite matching, task assignment with capacities, cut separation, and path packing.

\section*{Core Idea}

Dinic alternates between two steps:

\begin{itemize}
  \item run BFS from the source to build a level graph,
  \item run DFS only along edges that go one level deeper, pushing a blocking flow.
\end{itemize}

The level graph ensures that DFS only explores shortest augmenting paths in the current residual network.

\section*{Key Insight}

One BFS groups many augmenting paths together. Instead of finding and using only one path, Dinic pushes a whole
blocking flow before rebuilding levels. That is why it is much faster than plain Ford-Fulkerson with arbitrary paths.

\section*{Operations / Main Technique}

\begin{itemize}
  \item add a directed edge plus its residual reverse edge,
  \item BFS to compute levels,
  \item DFS with current-edge pointers to avoid rescanning dead edges,
  \item repeat until the sink becomes unreachable.
\end{itemize}

\paragraph{Worked example.}
In a bipartite matching reduction, the BFS builds layers \(s \to\) left part \(\to\) right part \(\to t\). The DFS
then pushes multiple augmenting paths through that layered structure before the next BFS is needed.

\section*{Correctness Intuition}

Every DFS push follows a valid residual path and preserves flow conservation. When the blocking flow step ends, every
source-to-sink path in the current level graph contains a saturated edge, so no more shortest augmenting paths remain.
If BFS later cannot reach the sink, the residual network has no augmenting path at all, so the current flow is maximum.

\section*{Complexity Analysis}

The general bound is \(O(V^2 E)\), with much better behavior on many contest graph families. In practice Dinic is the
standard max-flow baseline for sparse graphs.

\section*{Implementation}

The code uses:

\begin{itemize}
  \item adjacency lists of residual edges,
  \item level BFS,
  \item DFS with iterator pointers,
  \item 64-bit capacities.
\end{itemize}

\section*{Common Pitfalls}

\begin{itemize}
  \item Forgetting the reverse edge or updating the wrong reverse index.
  \item Using 32-bit capacities when the modeled totals are large.
  \item Omitting the current-edge pointer and then rescanning too much.
  \item Building the graph correctly but forgetting vertex-splitting when vertex capacities matter.
\end{itemize}

\section*{Variants / Extensions}

\begin{itemize}
  \item Bipartite matching as a unit-capacity flow problem.
  \item Lower-bound flow and circulation reductions.
  \item Dinic on graphs with vertex capacities via node splitting.
\end{itemize}

\section*{Practice Problems}

\begin{itemize}
  \item Maximum bipartite matching by max flow.
  \item Edge-disjoint paths between two vertices.
  \item Minimum cut partition after computing a max flow.
\end{itemize}

\section*{References}

\begin{itemize}
  \item \href{https://cp-algorithms.com/graph/dinic.html}{cp-algorithms: Maximum Flow - Dinic's Algorithm}.
  \item \href{https://usaco.guide/adv/max-flow}{USACO Guide: Max Flow}.
  \item \href{https://codeforces.com/catalog}{Codeforces Catalog}.
\end{itemize}
